\begin{alist}
\item Η γραφική παράσταση της συνάρτησης $f$ διέρχεται από τα σημεία $\text{A}(1,3)$ και $\text{B}(2,13)$, οπότε:
\[\begin{cases} f(1) = 3 \\ f(2) = 13 \end{cases} \Leftrightarrow \begin{cases} \alpha\cdot 2^1 + \beta = 3 \\ \alpha\cdot 2^2 + \beta = 13 \end{cases} \Leftrightarrow \begin{cases} 2\alpha + \beta = 3 \\ 4\alpha + \beta = 13 \end{cases} \overset{(-)}{\iff} \begin{cases} 2\alpha = 10 \\ 2\alpha + \beta = 3 \end{cases} \Leftrightarrow \begin{cases} \alpha = 5 \\ \beta = -7 \end{cases}.\]

\item Για $\alpha = 5$ και $\beta = -7$ η συνάρτηση γίνεται $f(x) = 5\cdot 2^x - 7$ και για $x = 0$ έχουμε $f(0) = 5\cdot 2^0 - 7 = 5\cdot 1 - 7 = -2$. \\
Άρα η γραφική παράσταση της συνάρτησης $f$ τέμνει τον άξονα $y'y$ στο σημείο $(0,-2)$.

\item Παρατηρούμε ότι για $x_1, x_2 \in \mathbb{R}$ με $x_1 < x_2$, έχουμε
\begin{align*}
2^{x_1} &< 2^{x_2} \Leftrightarrow \\
5\cdot 2^{x_1} &< 5\cdot 2^{x_2} \Leftrightarrow \\
5\cdot 2^{x_1} - 7 &< 5\cdot 2^{x_2} - 7 \Leftrightarrow \\
f(x_1) &< f(x_2).
\end{align*}
Άρα η $f(x) = 5\cdot 2^x - 7$ είναι γνησίως αύξουσα συνάρτηση.

\item Έχουμε
\begin{align*}
f(x) &> 4^x - 3 \Leftrightarrow \\
5\cdot 2^x - 7 &> (2^x)^2 - 3 \overset{2^x=y}{\iff} \\
y^2 - 5y + 4 &< 0 \quad (1).
\end{align*}
Το τριώνυμο $y^2 - 5y + 4$ έχει ρίζες $y_1 = 4$, $y_2 = 1$ (διότι $1+4=5=S$ και $1\cdot 4=4=P$). Η ανίσωση $(1)$ αληθεύει για $1 < y < 4$, δηλαδή:
\[1 < 2^x < 4 \Leftrightarrow 2^0 < 2^x < 2^2 \Leftrightarrow 0 < x < 2.\]
Τελικά η ανίσωση $f(x) > 4^x - 3$ αληθεύει για $x \in (0,2)$.
\end{alist}
